\subsection{Four unit-root ordinary-eighth families}
\label{proof:Q5-C096-ordinary}
After maximal genuine equal pairing let the residual lengths be
\begin{equation}\tag{Q96.0.1}\label{eq:Q96:0.1}
 \Sigma=(1,b_1,c_1,\ldots,b_k,c_k),\qquad
 1<b_1<c_1<\cdots<b_k<c_k,\quad k\ge1 .
\end{equation}
All \(L\ge1\) connector edges point from terminal to receiver,
whose existing unit arm is outward. For genuine equal-pair
products at the two roots, including empty products, put
\begin{equation}\tag{Q96.0.2}\label{eq:Q96:0.2}
 H=E_rpT^L(E_tF).
\end{equation}
The following directions are fixed by the residual lengths:
\[
\begin{array}{c|c|c}
\text{range}&F&\E H^8\text{ upper}\\\hline
b_1\ge4,\ k\ge1&q\prod_i h_{b_i}f_{c_i}&e^{-d}\\
b_1=3,\ k\ge2&qf_3h_{c_1}\prod_{i\ge2}h_{b_i}f_{c_i}&e^{-d/2}\\
b_1=2,\ k\ge3&qf_2h_{c_1}\prod_{i\ge2}h_{b_i}f_{c_i}&e^{-d/2}\\
\Sigma=(1,3,4)&qf_3h_4&1-d/601 .
\end{array}
\]
The first row retains the stronger all-domain rates
\begin{equation}\tag{Q96.0.3}\label{eq:Q96:0.3}
\begin{array}{ll}
b_1=4,k=1:&\E H^8\le e^{-2d},\\
b_1\ge4,k\ge2:&\E H^8\le e^{-[18+128(k-2)/5]d},\\
b_1\ge5,k\ge1:&\E H^8\le e^{-[3+128(k-1)/5]d}.
\end{array}
\end{equation}
Every instance of the first row keeps its displayed \(h_bf_c\) direction
even at \(L=1\). In each of the other rows the first pair has direction
\(f_bh_c\).

\begin{proof}
Apply Theorem~\ref{thm:unit-p:catalogue} with the actual terminal
minimum \(a=1\). Lemma~\ref{lem:unit-p:minimum-four} and
Section~\ref{part:unit-p:minimum-four-payments} give the three
strong rates in \eqref{eq:Q96:0.3}. Section~\ref{part:unit-p:three-later} gives the second row for every
finite \(L\), on the intervals meeting at \(d=1/6\).
Section~\ref{part:unit-p:two-many} gives the third row in the
same way. The fourth row is the genuine \(G_1\) case of
Section~\ref{part:unit-p:three-single}; its low-domain
finite continuation and the same-direction high comparison
meet at \(d=1/8\).

The original measure, terminal factor \(q\) and connector \(T^L\)
are retained throughout. Positivity restores \(E_t,E_r\);
complete reversal and \(d=0\) follow as in
Section~\ref{part:unit-p:decorations}.
\end{proof}
